% sec_1_3.tex — Sections 1–3 for arXiv submission
% Source: paper-draft-sec1-3.md (converted 2026-05-04)
% Includes §3.8 Pre-registration (not yet in main.tex monolith)
% Language: English (paper language)

% ---------------------------------------------------------------
% SECTION 1 — INTRODUCTION
% ---------------------------------------------------------------
\section{Introduction}

At 22:03 on April 25, 2026, an end-of-day cron job executed \texttt{nox-mem reindex}
without a dry-run flag. Within seconds, 183 entity records lost their \texttt{section},
\texttt{retention\_days}, and \texttt{section\_boost} fields---years of structured
operational context flattened into generic chunks. No error log. No alarm. The database
simply obeyed.

This incident did not expose a retrieval failure. It exposed a discipline failure. The
system had no enforced gate requiring a dry-run preview before a destructive reindex. It
had no mechanism for distinguishing between a trivial documentation update and a production
outage---both were stored with equal weight and retrieved with equal priority. And when six
AI agents (Atlas, Boris, Cipher, Forge, Lex, and Nox) shared a common operational corpus,
a single unchecked operation could silently degrade context for all of them simultaneously.

The incident became the motivating example for this paper.

\subsection{Problem Statement}

LLM agent memory is increasingly recognized as a prerequisite for sustained, coherent
assistance across sessions \cite{lewis2020rag, packer2023memgpt}. Yet the literature on
agent memory systems---from MemGPT's OS-inspired paging \cite{packer2023memgpt} to Mem0's
self-improving extraction \cite{chhikara2025mem0} to A-MEM's Zettelkasten-inspired linking
\cite{xu2025amem}---converges on a narrow set of concerns: retrieval accuracy on benchmark
corpora, latency at scale, and memory capacity. What these systems do not address is the
operational discipline required to keep a production memory system trustworthy over months
of continuous use.

Three specific gaps drive this work:

\textbf{Silent regression risk.} Ranking changes---new scoring formulas, boosting
coefficients, retrieval layer modifications---can degrade retrieval quality without any
observable error signal. Production memory systems have no established methodology for
validating such changes before they affect live queries. The April 25 incident is one
instance of a broader pattern: operations that silently alter system behavior with no guard
and no rollback.

\textbf{Incident severity as an invisible dimension.} Existing memory systems model recency
and semantic relevance. None, to our knowledge, model the operational cost of the
experience being stored. A lesson learned from a production outage and a routine
implementation note may share similar timestamps and embedding distances; they do not share
similar importance to future retrieval. This asymmetry is unrepresented in current
retrieval architectures.

\textbf{Multi-agent context silos.} When multiple agents operate over isolated memory
stores, intelligence accumulated by one agent is inaccessible to others. Solutions that
partition by \texttt{user\_id} \cite{chhikara2025mem0} or maintain per-agent state
\cite{packer2023memgpt} impose federation overhead or simply accept the silo. For trusted
multi-agent deployments---where agents serve a single operator with a shared context---this
isolation is an architectural choice, not a requirement.

\subsection{Thesis}

\textbf{Agent memory is not a retrieval problem---it is an operational discipline problem.}

Better embeddings, larger corpora, and faster fusion layers address the retrieval surface.
They do not address the question of whether the system can be trusted to evolve without
silent regressions, whether it surfaces the most operationally important context rather
than merely the most recent, or whether the intelligence accumulated by one agent is
available to the others who need it.

\subsection{Contributions}

We present \textbf{nox-mem}, a production memory system designed for multi-agent
operational use and built over approximately three months of continuous deployment.\footnote{Repository:
\url{https://github.com/totobusnello/memoria-nox}; submission tar at
\texttt{paper/publication/}; license MIT.} The system currently indexes 61,257 chunks
across a shared canonical corpus, maintains 99.97\% vector coverage, and operates at p95
latency below one second on commodity VPS hardware (8-core x86\_64) over the production
query mix. Three primary contributions distinguish this work:

\begin{enumerate}

\item \textbf{Pain-weighted salience as a secondary ranking modulator} (\S3.5). We
\textbf{propose} a retrieval scoring signal that explicitly models incident severity:
$\texttt{salience} = \texttt{recency} \times \texttt{pain} \times \texttt{importance}$,
where $\texttt{pain} \in [0.1, 1.0]$ encodes the operational cost of the recorded
experience---from trivial notes (0.1) to production outages (1.0). To our knowledge, this
is the first documented retrieval signal in the LLM agent memory systems literature that
treats incident severity as a first-class dimension; severity-aware prioritization itself
has well-established roots in observability and incident-response practice
(e.g., PagerDuty severity tiers, SIEM alert weighting), but has not been instrumented as
a retrieval ranking signal in the memory-systems literature surveyed in \S2.
Empirical ablation (\S5.4, E10, $n=31$
post-incident queries, hybrid mode) shows directional but non-significant aggregate effect
($\Delta = +0.0065$, 95\% CI $[-0.014, +0.034]$). The dimension provides meaningful lift
in the narrow regime where high-pain and low-pain chunks receive near-identical Gemini
semantic scores (Q55 case study, $\Delta = +0.349$). The transferable contribution is the
methodology of treating incident severity as typed schema input---a design pattern adoptable
independently of any retrieval stack---rather than a demonstrated corpus-wide retrieval
improvement.

\item \textbf{Enforced shadow discipline} (\S3.6). We describe a methodology for
validating ranking changes against a shadow telemetry baseline before deployment,
implemented as an architectural constraint rather than a documented suggestion. The
mechanism is codified via environment variable gating
(\texttt{NOX\_SALIENCE\_MODE=shadow|active}), enforced by a cron-checked health endpoint,
with a mandatory minimum observation period of seven days. During the salience validation
run (Fase~1.7b-b), this produced 191 promotion candidates, 16,608 review cases, and
45,743 archive recommendations before any activation. This is, to our knowledge, the first
memory system paper to codify such a discipline as an architectural constraint rather than
a suggested practice.

\item \textbf{Shared-canonical multi-agent design} (\S3.7). Six agents share a single
\texttt{chunks} table distinguished by \texttt{source\_file} prefix rather than separate
stores. Cross-agent retrieval requires no synchronization because partitioning never
occurred. When Forge records a decision about infrastructure, Atlas retrieves it directly
on the next relevant query---not because data was replicated, but because the design
assumed shared context from the start.

\end{enumerate}

Secondary contributions include: a reproducible evaluation harness with nDCG@10, MRR,
Recall@10, and Precision@5 over 50 curated golden queries; approximately three months of operational
evidence across five infrastructure upgrades without data loss; and comparative analysis
against strong retrieval baselines (see \S5).

\subsection{Paper Organization}

Section~2 surveys related work across knowledge graph--augmented retrieval, agent memory
systems, and production-oriented frameworks, identifying the specific gaps this work
addresses. Section~3 describes the system architecture across storage, retrieval, knowledge
graph, and operational discipline layers. Section~4 details the evaluation methodology,
including the golden query harness, baseline selection, and three-corpus design. Section~5
presents experimental results, including hybrid pipeline ablations, baseline comparisons,
pain dimension validation, and the shadow discipline case study. Section~6 discusses
deployment scope, limitations, and future work. Section~7 concludes.

% ---------------------------------------------------------------
% SECTION 2 — RELATED WORK
% ---------------------------------------------------------------
\section{Related Work}

Agent memory research has accelerated rapidly since the introduction of
retrieval-augmented generation \cite{lewis2020rag}, with recent work spanning knowledge
graph--augmented retrieval, autonomous memory management, and production deployment
frameworks. We survey the most directly relevant systems and identify the specific
dimensions where this work diverges.

\subsection{Knowledge Graph--Augmented Retrieval}

\textbf{GraphRAG} \cite{edge2024graphrag} introduced a pipeline that uses LLM extraction
to build entity-relation graphs over document corpora, applies community detection (Leiden
algorithm) to identify thematic clusters, and generates multi-level summaries for
query-focused retrieval. The approach substantially improves global query coverage---queries
requiring synthesis across many documents---compared to naive RAG. nox-mem adopts the
LLM-extraction paradigm for knowledge graph construction but diverges in objective:
GraphRAG optimizes for global summarization via community traversal, while nox-mem targets
entity-centric local retrieval for operational contexts, with closed-enum edge typing
designed to support downstream blast-radius analysis. GraphRAG does not model retrieval as
an operational concern and has no analogue to the discipline mechanisms described in
\S3.5--3.6.

\textbf{HiRAG} \cite{huang2025hirag} extends the KG-augmented paradigm with hierarchical
knowledge representation and iterative reasoning chains, demonstrating state-of-the-art
performance on multiple QA benchmarks. The hierarchical approach---traversing from coarse to
fine-grained graph levels---optimizes for accuracy on complex reasoning tasks. nox-mem
instead maintains a flat three-layer fusion pipeline (\S3.4), accepting the
accuracy-latency trade-off explicitly: p95 below one second at 61K chunks serves
operational use cases where query latency is a usability constraint, not just a benchmark
metric. HiRAG does not address multi-agent deployment, evaluation harnesses for operational
corpora, or retrieval discipline.

\textbf{HippoRAG} \cite{guo2024hipporag} draws on neurobiological memory models to improve
associative retrieval via hippocampal-inspired graph structures. The system demonstrates
that non-uniform memory consolidation---weighting some associations more strongly than
others based on structural graph properties---improves retrieval over flat indexing. nox-mem
introduces a conceptually related asymmetry through the pain dimension (\S3.5), though the
design rationale is grounded in operational incident management practice rather than
neuroscience. HippoRAG does not model incident severity as a retrieval signal.

\subsection{Agent Memory Systems}

\textbf{MemGPT} \cite{packer2023memgpt} proposed treating the LLM itself as an operating
system, with the agent autonomously managing its own context through function calls that
page between main context and external storage. The architecture gives agents direct
control over their memory---a compelling design for single-agent autonomy. nox-mem adopts
the opposite philosophy: memory is managed by an external pipeline (file watcher, ingest
router, nightly KG extraction), and agents are consumers of a shared canonical store rather
than managers of isolated state. This separation enables cross-agent retrieval as a
structural property while imposing the discipline that agents cannot corrupt their own
context.

\textbf{Mem0} \cite{chhikara2025mem0} introduced a production-focused memory layer with
LLM-driven memory extraction, deduplication, and temporal updates. Evaluated on the LOCOMO
benchmark, Mem0 demonstrates +26\% accuracy improvement over baselines on long-context
conversation scenarios. nox-mem shares the production emphasis but diverges in two
important respects. First, Mem0 relies on autonomous LLM editing decisions to update and
merge memories; nox-mem requires shadow-period validation before any retrieval-affecting
change propagates to production, specifically to mitigate silent regression risk. Second,
Mem0's multi-user design partitions memory by \texttt{user\_id}, treating agents as
user-addressable endpoints; nox-mem's shared canonical design treats agents as roles within
a single trusted operational context.

\textbf{A-MEM} \cite{xu2025amem} introduced agentic memory with Zettelkasten-inspired
structure: notes are automatically tagged, interlinked, and contextualized by an LLM
controller, enabling dynamic memory evolution. The approach handles unstructured
conversation well but does not address operational corpora with explicit schema
requirements, retention policies, or incident-severity differentiation. A-MEM has no
evaluation harness independent of LLM-generated quality judgments, and does not address
multi-agent deployments or ranking change discipline.

\subsection{Production-Oriented Frameworks}

\textbf{Cognee} \cite{topoteretes2024cognee} provides an open-source AI memory framework
implementing an Extract-Cognify-Load (ECL) pipeline with native knowledge graph support
and hybrid retrieval. Of the systems surveyed, Cognee is architecturally closest to
nox-mem: both maintain native KGs, both support hybrid retrieval, and both target
production deployment scenarios. Cognee achieves 3/5 on our architectural parity dimensions
(Table~\ref{tab:comparison}). The key differentiators are shadow discipline---Cognee provides
no enforced validation gate for retrieval-affecting changes---and the evaluation harness:
Cognee lacks a reproducible evaluation methodology with standard IR metrics over a fixed
golden set. nox-mem's eval harness (\S4) enables the kind of regression detection that the
shadow discipline gate is designed to act on.

\textbf{LangChain Memory} provides key-value and buffer memory primitives widely used in
production applications. The design prioritizes API simplicity and framework compatibility
over retrieval quality; it does not support hybrid retrieval, knowledge graph construction,
or evaluation. It serves as a practical lower bound in comparative analysis.

\subsection{Foundational Retrieval Methods}

\textbf{Retrieval-Augmented Generation} \cite{lewis2020rag} established the foundational
architecture combining parametric LLM knowledge with non-parametric document retrieval,
enabling factual question answering over dynamic corpora. nox-mem extends this foundation
to persistent, evolving operational corpora with explicit retention, pain weighting, and
multi-agent routing.

\textbf{Reciprocal Rank Fusion} \cite{cormack2009rrf} provides the rank-combination
mechanism at the core of nox-mem's hybrid retrieval layer. Cormack et al.\ demonstrated
that RRF consistently outperforms individual rank learning methods and Condorcet-based
fusion without requiring system-specific tuning. We adopt RRF with $k=60$ as the fusion
parameter, consistent with the paper's recommendations, as the final combination step over
BM25 and dense retrieval ranked lists.

\subsection{Positioning: Where nox-mem Fits}
\label{sec:positioning}

Table~\ref{tab:comparison} maps seven architectural and operational axes across the
surveyed systems. The table is intended to illustrate where the systems differ
structurally, not to function as a competition score. We include two dimensions where
nox-mem does not yet have coverage---corpus scale above 100K chunks and third-party
benchmark evaluation---because honest comparison requires acknowledging these gaps
explicitly. Limitations arising from these gaps are addressed in \S6.3; planned mitigations
are in \S6.5.

\begin{table}[!ht]
  \caption{Architectural comparison across memory systems. Axes reflect structural design
    choices, not performance rankings. \checkmark\ = implemented; $\approx$ = partial or
    indirect; $\times$ = absent.}
  \label{tab:comparison}
  \centering
  \small
  \resizebox{\textwidth}{!}{%
\begin{tabular}{llllll}
    \toprule
    System & KG Native & Hybrid Retr. & Eval Harness & Multi-Agent & Shadow Disc. \\
    \midrule
    \textbf{nox-mem (ours)} & closed-enum, 7 & FTS5+Gemini+RRF & nDCG/MRR/Recall & shared canonical & enforced $\geq$7d \\
    GraphRAG \cite{edge2024graphrag} & community det. & via KG queries & none & none & none \\
    MemGPT \cite{packer2023memgpt} & none & embedding-first & none & per-agent state & none \\
    Mem0 \cite{chhikara2025mem0} & optional (v2) & vector-only & LOCOMO only & user\_id part. & none \\
    A-MEM \cite{xu2025amem} & Zettelkasten & semantic-first & none & none & none \\
    HiRAG \cite{huang2025hirag} & hierarchical & multi-level & task-specific & none & none \\
    Cognee \cite{topoteretes2024cognee} & ECL pipeline & hybrid & ad-hoc & optional & none \\
    LangChain Memory & none & key-value & none & session\_id & none \\
    \bottomrule
  \end{tabular}}
  \begin{minipage}{\linewidth}
    \vspace{4pt}\footnotesize
    Five architectural axes are displayed; two additional axes---\emph{corpus scale
    \textgreater100K chunks} and \emph{3rd-party benchmark evaluation} (e.g., LOCOMO,
    LongMemEval, BEIR)---are summarized in the prose below. nox-mem currently covers
    5 of 7 axes ($\times$ on both omitted axes; see \S5.2 for partial BEIR/LOCOMO
    progress). Mem0 covers axis 7 via LOCOMO; no surveyed system has published a
    \textgreater100K-chunk evaluation. This is a comparison across architectural
    \emph{axes of difference}, not a competition score.
  \end{minipage}
\end{table}

We compare across seven axes. nox-mem covers 5/7: knowledge graph, hybrid retrieval, eval
harness, multi-agent design, and enforced shadow discipline. It does \textbf{not} yet cover
corpus scale above 100K chunks (current corpus: 61K) or evaluation against third-party
benchmarks such as LOCOMO, LongMemEval, or BEIR---the evaluation harness uses an internally
curated golden set. These are explicit limitations, not omissions, and are documented in
\S6.3 and \S6.5. The two axes with zero coverage across all surveyed systems---pain-weighted
salience and shadow discipline---are the primary novelty claims of this paper. Shadow
validation as a technique has well-established roots in production search evaluation
\cite{chapelle2012interleaved} and online controlled experiment methodology
\cite{kohavi2020trustworthy}. Our contribution is the instantiation of this discipline as
an enforced architectural constraint in an LLM agent memory system, rather than as a
testing methodology applied to web-scale traffic populations. Section~3 describes the
nox-mem architecture across all five covered dimensions; Sections~4--5 provide empirical
grounding.

% ---------------------------------------------------------------
% SECTION 3 — SYSTEM ARCHITECTURE
% ---------------------------------------------------------------
\section{System Architecture}

nox-mem is a production memory system built on SQLite with three primary subsystems---storage,
retrieval, and knowledge graph---plus an operational discipline layer that governs how the
system evolves. The implementation is in TypeScript using \texttt{better-sqlite3} for
synchronous database access, \texttt{sqlite-vec} for approximate nearest-neighbor search,
and Gemini for embedding generation and LLM-based KG extraction. The system has operated
continuously since March 2026 with 12 schema migrations and no irrecoverable data loss (one
metadata-loss event on 2026-04-25 was recovered via re-ingestion; see \S6.2).

\subsection{System Overview}

\begin{figure}[!ht]
  \centering
  \includegraphics[width=\linewidth]{figures/figure1-system-overview.pdf}
  \caption{System overview. Six personas read and write through a unified interface layer
    (CLI / MCP / HTTP) into a single canonical SQLite corpus (no per-agent silos).
    Retrieval runs as a three-layer hybrid (BM25 $\to$ 3072-d Gemini cosine $\to$ RRF,
    $k=60$). The shadow-mode pipeline and append-only \texttt{ops\_audit} are cross-cutting
    governance: every ranking change ships shadow-first; every destructive op is wrapped by
    \texttt{withOpAudit()} snapshot.}
  \label{fig:overview}
\end{figure}
\FloatBarrier

The entry point for all CLI operations is \texttt{dist/index.js} (the compiled package
binary). An HTTP API on port 18802 exposes search, knowledge graph, health, and operational
endpoints for agent integration.

\medskip
\begin{verbatim}
Source files (.md / .docx / .pptx / .pdf / entity files)
    |
    v
ingest-router (inotifywait watcher + manual CLI)
    |
    +------> chunks table (canonical, 61K rows)
                |
                +------> chunks_fts (FTS5 BM25 index)
                +------> vec_chunks (sqlite-vec, 3072d embeddings)
                +------> kg_entities + kg_relations (LLM-extracted, nightly)
                +------> search_telemetry (shadow-mode + eval harness)
                |
    query --> search() pipeline:
                Layer 1: FTS5 BM25 candidates
                Layer 2: Gemini semantic candidates
                Layer 3: RRF fusion (k=60) + salience boost + section_boost
                --> ranked results (p95 < 1s @ 61K chunks)
\end{verbatim}
\medskip

\subsection{Storage Layer}

The canonical store is a single SQLite database accessed via \texttt{better-sqlite3} for
synchronous, blocking reads---a deliberate choice that simplifies the concurrency model for
a single-node deployment. Three coordinated tables implement the storage layer.

\texttt{chunks} is the central table, containing content text, metadata
(\texttt{source\_file}, \texttt{chunk\_type}, \texttt{created\_at}, \texttt{last\_seen}),
and operational fields (\texttt{pain} REAL DEFAULT 0.2, \texttt{importance} REAL,
\texttt{section} TEXT, \texttt{section\_boost} REAL, \texttt{retention\_days}). The schema
has evolved through 12 versions (v1 through v12) since March 2026; all migrations have been
applied to the live database without downtime or data loss. Key schema additions include
\texttt{retention\_days} (v8, typed retention by chunk category), \texttt{pain} (v9,
severity signal), and \texttt{section} with \texttt{section\_boost} (v10, structural
weighting for entity file sections).

\texttt{chunks\_fts} is an FTS5 virtual table providing BM25-ranked full-text search over
chunk content. FTS5 vanilla AND-mode behavior---requiring all query terms to
match---means that natural language queries frequently return zero results when run against
the lexical index alone. This structural characteristic of FTS5 makes hybrid retrieval a
necessity rather than an optimization, a finding confirmed in evaluation (\S5.1).

\texttt{vec\_chunks} stores 3072-dimensional Gemini embeddings
(\texttt{gemini-embedding-001}) via the \texttt{sqlite-vec} extension, enabling
approximate cosine similarity search without external vector store infrastructure. Current
vector coverage is 99.97\% across 61,257 chunks.

Retention policy is encoded directly in the schema via \texttt{retention\_days}: feedback
and person records are set to \texttt{NULL} (never decay); decisions and projects have
365-day retention; team context 120 days; daily notes 90 days; pending items 30 days. This
ensures that the system's own operating rules cannot be silently evicted by a generic decay
policy.

\subsection{Retrieval Layer}

The retrieval pipeline operates in three sequential layers, with results from layers one and
two merged via Reciprocal Rank Fusion \cite{cormack2009rrf}:

\textbf{Layer 1---Lexical retrieval.} FTS5 BM25 search expands the query against the
\texttt{chunks\_fts} index, returning up to $k_1$ ranked candidates. Results are
deduplicated by chunk ID before passing to fusion.

\textbf{Layer 2---Semantic retrieval.} The query is embedded using
\texttt{gemini-embedding-001} (3072-d), and cosine similarity search is performed against
\texttt{vec\_chunks} via \texttt{sqlite-vec}, returning up to $k_2$ ranked candidates.

\textbf{Layer 3---RRF fusion and boosting.} Candidates from layers one and two are merged
via Reciprocal Rank Fusion \cite{cormack2009rrf} with $k=60$:

\[
  s_{\text{rrf}}(d) = \frac{1}{k + r_{\text{fts}}(d)} + \frac{1}{k + r_{\text{sem}}(d)}
\]

where $r_{\text{fts}}(d)$ and $r_{\text{sem}}(d)$ are the BM25 and cosine-similarity ranks
of chunk $d$ respectively.

The fused score is then combined with two boost signals via log-additive aggregation:

\[
  \log s_{\text{final}}(d) = \log s_{\text{rrf}}(d) + \log s_{\text{salience}}(d)
    + \log s_{\text{section}}(d)
\]

where $s_{\text{salience}}(d) = \texttt{recency}(d) \times \texttt{pain}(d) \times
\texttt{importance}(d) \in [0, 2]$ (\S3.5), and $s_{\text{section}}(d) \in
\{2.0,\,1.5,\,0.8,\,1.0\}$ for compiled, frontmatter, timeline, and unstructured (NULL)
section types respectively.

Equivalently in linear space,
$s_{\text{final}}(d) = s_{\text{rrf}}(d) \cdot s_{\text{salience}}(d) \cdot s_{\text{section}}(d)$.
The two formulations are mathematically identical; the log-additive form is operationally
preferred because it allows each boost's contribution to be recorded and monitored
independently in the \texttt{search\_telemetry} log without re-multiplication.

Top-$k$ chunks are returned in descending order of $s_{\text{final}}$.

The pipeline achieves p95 search latency below one second across the 61K chunk corpus,
measured in production over the observation period. The \texttt{--no-hybrid} flag disables
semantic retrieval for latency-constrained contexts where lexical precision is acceptable.

\subsection{Knowledge Graph Layer}

An LLM-extracted knowledge graph augments chunk retrieval with structured relationship
data. Entities and relations are extracted incrementally via nightly batch runs using
\texttt{gemini-2.5-flash} and stored in \texttt{kg\_entities} ($\approx$402 entities) and
\texttt{kg\_relations} ($\approx$544 relations).

Entity types follow a closed taxonomy of 10 categories: person, project, agent, tool,
concept, organization, technology, document, decision, and metric. Relation edges use
free-form \texttt{relation\_type} text combined with a closed-enum
\texttt{relation\_reason} field with seven canonical values: \texttt{depends\_on},
\texttt{derived\_from}, \texttt{opposes}, \texttt{extends}, \texttt{replaces},
\texttt{mentions}, and \texttt{unknown}. The closed enum enables downstream
tooling---\texttt{impact~\textless entity\textgreater} computes one-hop blast radius with
reason-priority weighting; \texttt{detect-changes} maps git diffs to affected KG
entities---without requiring LLM inference at query time.

Edge type enum coverage required explicit engineering. Initial extraction with an optional
enum field and a ``use unknown if unsure'' prompt instruction produced 86\% \texttt{unknown}
labels across n=100 sampled relations. A three-path fix---a revised prompt with explicit
examples for each non-unknown category; a code-side defensive normalization map
(\texttt{RELATION\_TYPE\_TO\_REASON}, 24 input aliases covering PT-BR and EN free-text
variants that collapse to the 7 canonical reasons); and a post-extraction validation
pass---improved the enum coverage rate from 14\% to 56\% ($4\times$ improvement),
equivalently reducing the \texttt{unknown} rate from 86\% to 44\% (n=100 sample). The
enum coverage rate measures the proportion of sampled relations where the LLM committed to
one of the six non-unknown typed values rather than falling back to \texttt{unknown}; it is
a self-reported coverage rate, not a classification accuracy against a human-annotated
ground truth (see \S6.3). This experience reinforces a general principle: LLM-extracted
structured fields with optional enums and uncertainty-deferring prompts will systematically
underperform; defensive code-side normalization is required.

\subsection{Pain-Weighted Salience}

The salience formula is:

\begin{equation}
  \texttt{salience}(\mathit{chunk}) =
    \texttt{recency}(\mathit{chunk})
    \times \texttt{pain}(\mathit{chunk})
    \times \texttt{importance}(\mathit{chunk})
  \label{eq:salience}
\end{equation}

\noindent where:
\begin{itemize}
  \item $\texttt{recency} \in [0,\,1]$ is an exponential decay function over the
    \texttt{last\_seen} timestamp, parameterized by the chunk's \texttt{retention\_days}
    value;
  \item $\texttt{pain} \in [0.1,\,1.0]$ encodes the operational severity of the recorded
    experience, annotated via \texttt{pain: 0.X} markers in entity files. Values are
    assigned on a five-point semantic scale: 0.1 (trivial note), 0.3 (minor friction),
    0.5 (significant issue), 0.7 (major incident), 1.0 (production outage with data loss
    or downtime);
  \item $\texttt{importance} \in [0,\,1]$ is derived from \texttt{mention\_count} across
    the corpus and entity-type priors.
\end{itemize}

The formula is multiplicative: a high-pain chunk with low recency can still score above a
low-pain chunk with high recency when the severity differential is large. This reflects an
engineering choice motivated by operational practice rather than a biological or
psychometric claim. The five-point scale follows the structure of established incident
severity taxonomies \cite{pagerduty2023severity,beyer2016site}: just as incident management
systems assign higher dispatch priority to a production outage than a documentation note
regardless of their timestamps, pain weighting ensures that incident-derived memory
dominates retrieval over routine content even when recency is comparable.

The salience signal is exposed in shadow mode at \texttt{/api/health.salience} without
affecting production ranking, enabling comparison against the baseline before activation
(\S3.6). Ablation results comparing retrieval quality with real versus uniform pain values
on $n=31$ post-incident queries are reported in \S5.4 (E10, 2026-05-04); the aggregate
result is directional but not statistically significant, with Q55 as the primary positive
case study ($\Delta = +0.349$).

\FloatBarrier
\subsection{Enforced Shadow Discipline}

\begin{figure}[!ht]
  \centering
  \includegraphics[width=0.9\linewidth]{figures/figure3-shadow-state-machine.pdf}
  \caption{Shadow discipline state machine. Every ranking-affecting change traverses this
    machine. Activation requires three independent signals: at least seven days of shadow
    telemetry, a healthy distribution (no degenerate clustering), and explicit human
    approval. Regression at any point reverts via a single environment-variable
    flip---zero schema rollback required.}
  \label{fig:shadow}
\end{figure}
\FloatBarrier

The shadow discipline mechanism enforces a separation between observing a retrieval change
and deploying it. Any modification that affects ranking---scoring formula changes, boost
adjustments, new retrieval layers---must run in shadow mode for a minimum of seven days
before activation. This is an architectural constraint, not a documented recommendation.
Shadow validation has well-established roots in production search and online controlled
experiments \cite{kohavi2020trustworthy,chapelle2012interleaved}. Our contribution is the
application of shadow discipline specifically to LLM agent memory systems, with
architectural enforcement via cron and health endpoints.

Implementation relies on three components: (1)~an environment variable gate
(\texttt{NOX\_SALIENCE\_MODE=shadow|active}) that controls whether the salience score
affects production ranking or is only written to \texttt{search\_telemetry}; (2)~a health
endpoint (\texttt{/api/health.opsAudit}) that cron checks every 15 minutes and alerts via
Discord if the shadow observation period has not been met; (3)~an append-only audit log
(\texttt{ops\_audit} table) with database triggers that block DELETE and UPDATE on rows
with terminal status, ensuring the shadow-period record cannot be retroactively modified.

The distinction from A/B testing is important: A/B testing requires live traffic over a
production population. Shadow discipline for operational/personal corpora has no such
traffic---queries are infrequent, irregular, and non-i.i.d. The shadow mechanism instead
collects telemetry on what the proposed ranking \emph{would have done} on each live query,
accumulating promotion, review, and archive candidate distributions across the observation
period. The distribution comparison then drives the activation decision.

The salience validation run (Fase~1.7b-b) ran for seven days before activation, producing
191 promotion candidates, 16,608 review cases, and 45,743 archive recommendations. The
distribution was reviewed, found consistent with design intent, and activation proceeded.
Counterfactual analysis suggests that the April~25 reindex incident---the motivating example
in \S1---would have been detected during a shadow period: the reindex operation altered
chunk distributions in ways that would have surfaced as anomalous in telemetry within hours.

\subsection{Shared-Canonical Multi-Agent Design}

Six agents (Atlas, Boris, Cipher, Forge, Lex, and Nox) share a single \texttt{chunks}
table. Agent-specific content is distinguished by \texttt{source\_file} path prefix
(\texttt{agents/<name>/...}) rather than separate tables, schemas, or databases.
Cross-agent queries use \texttt{cross-search} and \texttt{cross-kg} CLI commands, which
apply additional \texttt{source\_file} filters to retrieve context generated by a specific
agent.

The design assumes a single trusted operator. It is explicitly not suitable for
multi-tenant SaaS deployments, where user isolation is a security requirement rather than
an engineering choice. Within the trusted-operator constraint, shared-canonical design
provides cross-agent intelligence as a structural property: when Forge ingests a decision
about an infrastructure change, Atlas retrieves it on the next architecturally relevant
query without any synchronization step.

The operational tooling layer---\texttt{detect-changes}, \texttt{impact},
\texttt{api-impact}, \texttt{consolidate-merge}, \texttt{kg-reclassify}---extends the
system to support multi-agent workflows beyond retrieval, enabling agents to query the KG
for blast-radius analysis before executing changes, and to surface merge candidates across
agent-sourced entities.

The architecture in \S3 provides the substrate for the empirical evaluation in \S4--5.
Section~4 describes the evaluation methodology, including the golden query harness,
baseline systems, and three-corpus protocol designed to address the single-corpus and
single-curator limitations of operational memory benchmarking.

\FloatBarrier
\subsection{Pre-registration and Reproducibility Commitments}

A persistent concern in machine learning evaluation is post-hoc query set construction:
a curated golden set assembled with knowledge of retrieval results can unintentionally
favor the system under evaluation, inflating reported metrics without any deliberate
manipulation \cite{munafosystematic2017}. We address this concern explicitly by anchoring
the experimental design to a pre-registration artifact that predates all result collection.

\textbf{Pre-registration document.} The evaluation protocol---including query categories,
metrics, baseline selection criteria, and success thresholds for each
hypothesis---was committed to \texttt{paper/publication/03-experiments-needed.md} before
any baseline or ablation experiment was executed. This document is tracked in the public
repository (\url{https://github.com/totobusnello/memoria-nox}) and its commit history is
immutable. The git commit hash and ISO-8601 author timestamp for that file should be
verified at submission via:

\medskip
\begin{verbatim}
git log --follow -1 --format="%H  %aI" \
    paper/publication/03-experiments-needed.md
\end{verbatim}
\medskip

\textbf{Golden query set freeze.} The 60-query golden set (\texttt{eval/golden-queries.jsonl},
comprising the 50-query main set R01b and the 10-query held-out subset R01c) is included
in the public repository and the submission tar. The verifiable identifiers are:

\begin{itemize}
  \item \textbf{Git commit hash:} \texttt{f75d1864b581d0a0933704e997669bbd41aa507f}
  \item \textbf{Author timestamp (ISO-8601, BRT):} \texttt{2026-05-04T13:38:01-03:00}
  \item \textbf{Freeze mtime (VPS, UTC):} \texttt{2026-05-04T09:45Z}
  \item \textbf{SHA-256:} \texttt{9bff8ee7b9056eff6a1af22305cae762aa4b98e682578faffb4c22cdd0a2cd7d}
  \item \textbf{File size:} 8990 bytes, 60 lines (one query per line, JSONL)
\end{itemize}

The SHA-256 hash allows any third party to verify byte-for-byte that the included
\texttt{golden-queries.jsonl} matches the artifact used for all reported evaluations.
The golden query set was frozen \emph{before} any BM25 (Pyserini), multilingual-e5-base,
BEIR TREC-COVID, LOCOMO, or pain-ablation (E10) experiment was executed; all results in
\S5 were produced against this exact file.

\textbf{Held-out subset disclosure.} The 10 queries designated R01c (held-out for baseline
robustness validation) were locked as a named subset within the golden set before the final
tuning sprint that fixed the RRF $k=60$ parameter and the salience boost coefficients.
They were not constructed after tuning; they are a partition of the original 60-query set,
identified by their position index (Q51--Q60) in the JSONL file. R01c is evaluated at most
once per major system revision and its results are reported separately from the 50-query
main set to prevent iterative query refinement bias (\S4.1).

\textbf{Registered reports context.} While this paper does not constitute a formal
registered report in the sense of \cite{munafosystematic2017}---no journal or conference
pre-acceptance was obtained prior to experimentation---the \texttt{03-experiments-needed.md}
document serves an equivalent function within the project's operational discipline: it
records pre-registered hypotheses with explicit success thresholds (e.g.,
$\Delta$~nDCG@10 $\geq 0.05$ for E10) that are reported honestly regardless of outcome,
as evidenced by the DIRECTIONAL, NOT SIGNIFICANT result for the E10 pain ablation (\S5.4).
The pre-registration document, the golden query set, and all evaluation scripts are
included in the public repository under MIT license, permitting full independent
verification.
